% Executor candidate theorem. H3 remains OPEN because C_mu is not discharged.
\begin{theorem}[Conditional reachable exponent/BerExp bridge]\label{thm:H3-conditional}
Fix a scalar call made by the committed FT1536 ternary signer after an H4
accepted tree has been decoded.  Let $\mu$ be its binary64 center and impose
the explicitly named integer-safety premise
\[
 \mathsf C_\mu:\qquad
 -2147483283\le \lfloor\mu\rfloor\le2147483281.
\]
Write $s=\lfloor\mu\rfloor$ and let
\[
 \rho=\mu-s\in[0,1),\qquad
 \widehat r=\operatorname{RN}(\rho)\in[0,1]
\]
be respectively the exact real fractional part and the committed binary64
subtraction result.  Then the source \texttt{fpr\_floor} result converts to
\texttt{int s} safely, \texttt{fpr\_of(s)} is exact, and
\[
 |\widehat r-\rho|\le 2^{-53}+2^{-1075}.
\]
Every current proposal return $s+z$, for $-365\le z\le366$, is a defined
\texttt{int} result.  The endpoint $\widehat r=1$ is allowed.

Let $d$ be the exact dyadic represented by the computed binary64
\texttt{dss}, let $d_0$ be the first eligible coefficient in the committed
five-bank selector, and let $k$ be the table output.  Put
\[
 \delta_0=\rho,\qquad \delta_1=1-\rho,\qquad
 X=k^2(d-d_0)+(2k\delta_b+\delta_b^2)d.
\]
For the finite implemented proposal support, let $K_{\rm prod}$ use the
actual source values
\[
 \widehat\delta_0=\widehat r,\qquad
 \widehat\delta_1=\operatorname{RN}(1-\widehat r),
\]
the production binary64 exponent chain, range reduction, evaluator, prefix
and comparator.  Let $K_{\rm exact}$ accept the same proposal with exact
probability $e^{-X}$, so that its rejection identity is centered at the
original target $s+\rho=\mu$.  Uniformly over the complete H4 width interval,
both scalar branches, all five selector domains, all table-reachable $k$, and
all source-reachable binary64 centers satisfying $\mathsf C_\mu$,
\[
 \lVert\widehat\pi K_{\rm prod}
       -\widehat\pi K_{\rm exact}\rVert_1
 \le D_{\rm H3},\qquad
 \log_2D_{\rm H3}<-44.9362260.
\]
The bound includes the $s\ge64$ source cutoff.  Pointwise on that region the
discarded exact probability is less than $2^{-63.9999999999998}$.  On the
main region the complete relative likelihood defect has logarithm below
$-44.9362286$.

If the exact-kernel acceptance floor $a=9/20$ from H1R is used, then the
normalized accepted scalar laws obey
\[
 \operatorname{TV}\bigl((\widehat\pi K_{\rm prod})_{\rm acc},
                         (\widehat\pi K_{\rm exact})_{\rm acc}\bigr)
 \le {D_{\rm H3}\over a-D_{\rm H3}},
 \qquad
 \log_2\!\left({D_{\rm H3}\over a-D_{\rm H3}}\right)<-43.7842229.
\]
Consequently sequential maximal coupling gives, conditionally on
$\mathsf C_\mu$ at every adaptive call, logarithmic TV bounds
$-32.1992604$ for the 3072 outputs of one \texttt{do\_sign} trial and
$-28.1992604$ for at most 16 trials with an explicit $\bot$ output.
\end{theorem}

\begin{proof}
The displayed integer interval and parsed table support make the two source
integer operations safe.  Correct RN-even subtraction gives
$|\widehat r-\rho|\le u\rho+\eta\le u+\eta$, for $u=2^{-53}$ and
$\eta=2^{-1075}$.  The minimum negative subnormal shows that this subtraction
need not be exact: $\mu=-2^{-1074}$ gives $s=-1$, exact
$\rho=1-2^{-1074}$ and computed $\widehat r=1$.  For branch zero the delta
error is at most $u+\eta$; for branch one the additional RN-even subtraction
gives at most $2(u+\eta)$.  Since all four exact/computed deltas lie in
$[0,1]$, the resulting semantic exponent displacement is at most
\[
 2(u+\eta)(2k+2)d.
\]

The H4 endpoint replay, including the terminal $2/\sqrt3$ multiplication,
gives the exact computed range
\texttt{dss} in
\texttt{[0x3f455555c49559f4,0x3fd203d2c1ca3682]}.
The lower endpoint exceeds the final bank coefficient.  Direct parsing of
all 512 thresholds in each bank gives maxima
$k=(29,59,118,235,365)$.

At every finite RN-even operation
$|\operatorname{RN}(y)-y|\le u|y|+\eta$.  Propagating this inequality through
the seven displayed source operations, and adding the separate center-split
semantic displacement above, handles normal, subnormal and zero-rounded
terms without a universal relative-error assumption.  The complete domain
gives $\widetilde X$ below the first range-reduction band with $s=394$.  On
$s\le63$ the revised mixed chain bound has logarithm below $-45.2773722$.

An exact monotone binary64 partition covers every $\widetilde X$ in the 394
bands $s=0,\ldots,393$.  The stored \texttt{log(2)} and
\texttt{inv\_log(2)} constants leave nonnegative residuals, with a maximum
of \texttt{0x3fe62e42fefa3a00}, exactly 17 binary64 points above the stored
\texttt{fpr\_log2}.  H2 covers the real interval $[0,\ln2)$; direct RIF and
integer-recurrence evaluation covers those 17 fringe points at the same
$2^{-50}$ budget.  The main range-reduction displacement has logarithm below
$-47.6656431$.

For $s\le63$, the prefix succeeds with probability exactly $2^{-s}$ and the
55-bit comparison succeeds with probability exactly $z/2^{55}$.  Combining
the split, chain, range and H2 bounds gives the stated relative defect.  For
$s\ge64$ the source forces zero; the exact first-$s=64$ binary64 threshold
minus the global displacement gives the stated pointwise cutoff mass.
Summing absolute accepted-submeasure differences gives $D_{\rm H3}$.
If the exact mass is at least $a$, the production mass is at least $a-D$ and
normalization costs at most $D/(a-D)$.  Adaptive sequential coupling proves
the two multiplicity statements.
\end{proof}

\begin{remark}[Open integer-reachability premise]
The committed source checks only that $\mu$ is finite before assigning
\texttt{fpr\_floor(mu)} to an \texttt{int}; H4 constrains leaves and widths,
not internal LDL coefficients or centers.  The signer-accepted key population
does not require KeyGen's Babai representative: the exact shear
$(F,G)\mapsto(F+kf,G+kg)$ preserves NTRU and H4 while the root center changes
by an integer convolution.  Coarse source-bound decoder estimates already
exceed \texttt{INT\_MAX}; the canonical-KeyGen route lacks a uniform
Babai/FFT/LDL error certificate.  Therefore this theorem remains conditional
and H3 remains open.
\end{remark}

\begin{remark}[Future fail-closed route]
The separately planned task \texttt{S20-H3G-CENTER-RANGE-GATE-001} specifies
the half-open pre-floor check
$-2147483283\le\mu<2147483282$, a widened checked \texttt{s+z}, and sticky
fault propagation to $\bot$.  This executor bundle does not modify source.
For a future guarded theorem, gate failure must be common $\bot$ on both
sides and gate success invokes the theorem above.
\end{remark}

\begin{remark}[Comparator floor and E1/E2 boundaries]
The comparator probability grid has spacing $2^{-55}$, hence a uniform
continuous-target architecture cannot beat the half-step floor $2^{-56}$.
E1's refit and E2's U72 chain are not deployed.  E2's subnormal counterexample
is why the proof above is mixed absolute rather than universally relative.
\end{remark}
